\section{Background: the evolution of FishBot}
\label{sec:background}

\fishseven{} is the seventh configuration in a lineage, and several of its components and most of its
evaluation machinery exist because of what earlier versions found. This section is deliberately
short. Its purpose is to explain why \fishseven{} looks the way it does, not to re-report earlier
studies; the detailed corrections those studies require are collected in Appendix~\ref{app:corrections}.

\subsection{v0.2 to v0.4: engine, inference and policy foundations}
\label{sec:background-early}

The first versions established the domain and the inference. A count-conditioned, empirically tuned
linear policy in a scripting-language simulator became a rules-exact C++ engine carrying an exact
reference posterior, a faster deployed approximation, a twenty-feature ask score and a value-based
declaration rule. The durable contribution of that period is the inference object of
\S\ref{sec:inference}: an exactly computable observer-conditioned deal posterior, cheap enough to run
in all six seats at once.

That period also established a habit of evaluation this paper spends considerable effort undoing.
Cells were a few thousand games against a handful of opponents, ablations were unpaired, and
conclusions were drawn from point estimates.

\subsection{v0.5: a failure mode removed, and honestly reported as worth nothing}
\label{sec:background-vfive}

v0.5 diagnosed a two-question deadlock in mirror play, in which two policies alternately ask each
other for cards neither will surrender until the game ends in misdeclaration, and removed it. It then
reported that removing it was worth almost nothing in win rate.

That report is the honest ancestor of this one, and it is where the project's central difficulty
first became visible. A release that removes a genuine defect and gains no measurable strength is
either at the ceiling of its policy class or is measuring badly, and those two diagnoses call for
opposite responses.

\subsection{v0.6: stronger play, and a measurement crisis}
\label{sec:background-vsix}

The v0.6 cycle set out to recover several points of headroom that the accumulated design register
described. It found almost none of it, and discovered in the process that much of the measurement
machinery was unsound. Four results from that cycle bear directly on \fishseven{}.

\paragraph{Exchangeable ties are irreducible, so the tie-break is a coordination decision rather than
a strength decision.} At \vsixTieShareFive\% of \emph{contested} ask decisions two or more candidates
score bit-for-bit identically, and the winner is whichever the enumerator emitted first.
\vsixTieSameSuitFive\% of those ties are two cards of one half-suit at one target, and the exact
posterior \emph{separates} them in \vsixTieExactSepPctFive\% of cases, meaning it assigns them equal
probability in every case. Every rule that ranks tied candidates by that probability therefore
realises the same hit rate, and so does enumeration order. The largest item on the inherited design
register was not a defect but a property of the information set. What remains addressable is the
dependence on enumeration order itself, which is a correctness problem, and
\S\ref{sec:agent-tiebreak} is how \fishseven{} addresses it.

\paragraph{The exact posterior is a worse predictor than the deployed approximation.} Scored as
predictors on identical states, the ranking is:

\begin{center}\small
\begin{tabular}{lrr}
\toprule
Posterior & Log loss (nats) & Argmax hit rate \\
\midrule
Exact, uniform deal prior & \vsixBeliefExactNLL{} & \vsixBeliefExactHit\% \\
Rules posterior, no policy prior & \vsixBeliefNoPriorNLL{} & \vsixBeliefNoPriorHit\% \\
\textbf{Deployed approximation} & \textbf{\vsixBeliefShippedNLL{}} & \textbf{\vsixBeliefShippedHit\%} \\
\bottomrule
\end{tabular}
\end{center}

The exact object is the worst row. The policy prior accounts for the entire difference: exactness in
this game is purchased by discarding behaviour, because the block construction of
\S\ref{sec:inference} has no parameter through which an opponent's tendencies could enter. Exactness
and behaviour-awareness are not available in the same object here. This closed an experiment three
studies had listed as outstanding, and it established a methodological rule the project has applied
since, namely that exact Bayesian inference is an assumption to be tested rather than a property to
be assumed. \S\ref{sec:development-alloc} is the second occasion on which that rule paid.

\paragraph{Determinized search suffers an optimiser's curse, and a paired guard removes it.} v0.6
built a search that samples full deals from the public posterior and plays each forward with all six
players reconstructed at their own information sets, rather than made clairvoyant, which
\S\ref{sec:game-why} explains would be degenerate. Chosen by an unguarded argmax over sampled means,
that search is \vsixCurseGap{}~pp \emph{worse} than the blueprint it searches from, measured against a
control that forces the same code to return the blueprint's own choice at identical budget, seed and
sample size. The cause is the statistic rather than the search: one determinization's return is the
final half-suit differential, whose spread is an order of magnitude larger than genuine differences
between candidates, so the argmax selects on residual noise and does so more confidently the more
candidates it is offered. A paired lower-confidence-bound deviation rule recovers all of it. The
guarded search, which v0.6 measured and then deployed switched off, is a component of \fishseven{}
(\S\ref{sec:agent-search}).

\paragraph{The evaluation could not resolve a one-point effect, and the optimiser could not move a
policy.} Pooling every same-pair cell in the v0.6 battery gave \vsixHeadPooled\% with a naive $z$ of
\vsixHeadZ{}, a null that an earlier draft of that report published as one. Re-run at high power
across \vsixBigFiveBanks{} disjoint banks, the same comparison is \vsixBigFive\% over
\vsixBigFiveN{} games $[\vsixBigFiveLo, \vsixBigFiveHi]$, above parity on all \vsixBigFiveAbove{}
banks. The effect had always been present; the instrument could not see it. In the same cycle, five
mechanisms cleared a 95\% paired interval at the per-cell budget the project had used since v0.3 and
returned exactly zero at three times that budget. Separately, the fitting harness was found to be
incapable of moving a policy at all, using one step size for thirty-four coordinates spanning two and
a half orders of magnitude, an unpaired objective, and an objective documented as a worst case that
was in fact a weighted mean.

\subsection{What v0.6 handed to the v0.7 programme}
\label{sec:background-handoff}

A policy roughly nine tenths of a point stronger than its predecessor; a design register most of
whose entries had just been closed as worth nothing; a repaired optimiser; and an explicit diagnosis
that at the budget the project had always used, the policy class looked flat without it being known
whether the class \emph{is} flat or the instrument was still too coarse.

The v0.7 programme was organised to answer that question in a way that could not be argued into a
positive result. Five phases ran under separate briefs, each starting from a cleared context:
characterise the instrument and establish what it can detect; search for adversaries in the open
against the \vsix{} frontier rather than against the deployed policy; build and test candidate
mechanisms; select and freeze exactly one configuration; and evaluate that configuration on sealed
material under a protocol registered in advance. \S\ref{sec:development} reports the first three
phases and \S\ref{sec:results} the last.
